\documentclass[11pt,a4paper]{article}
\usepackage[utf8]{inputenc}
\usepackage[T1]{fontenc}
\usepackage[spanish]{babel}
\usepackage{hyperref}
\usepackage{geometry}
\geometry{margin=2.5cm}
\usepackage{parskip}

\title{352 no usar YO sin individualidad}
\author{Alejandro Valenzuela Rivera\\\small Universidad Autónoma de Nuevo León}
\date{\today}
\begin{document}
\maketitle

\begin{quote}
\textbf{Sesgo declarado.} Este trabajo se sostiene desde el sesgo declarado de su autor, y no finge neutralidad. Lo que afirma se puede tumbar, y ese es el punto: el falsificador está escrito, no escondido.
\end{quote}

\section*{Resumen}
El proceso no se escribe en primera persona sin declarar quién la tiene: un YO sin individualidad es una voz que se atribuye una mirada que no tiene.

\section{Por qué se sostiene}
Un proceso que dice 'yo' sin declarar la individualidad se vuelve autor sin serlo, y el error no tiene a quién reclamarle.

\section{De dónde viene}
La autoridad, en mayúsculas: 'procesos no mis creencias, por eso también la regla de no usar YO sin individualidad'.

\section{El caso}
El repo: todo proceso habla en tercera persona, y el autor firma con su nombre.

\section{Método}
\textquestiondown{}El 'yo' del texto tiene un tipo:nombre que lo respalde?

\section{Límites, y qué lo tumba}
Un proceso que se atribuye una vista sin declarar de quién es.

\section{Lo que queda abierto}
\textquestiondown{}El YO del proceso declara quién mira?

\section*{Cómo se leyó esta tesis}
Esto se escribió en el repositorio man-made-empathy, del que este candado es una pieza: candados/04 el lenguaje/352 no usar YO sin individualidad.md. La máquina que lo compuso lo hizo desde este mismo archivo, y no agregó nada que no estuviera escrito. Si algo de esto es falso, se dice en el mismo sitio, y se corrige con fecha.

\end{document}
